\documentclass{article}
\usepackage{amsmath,amssymb}
\usepackage[margin=0.65in]{geometry}
\allowdisplaybreaks
\usepackage[UTF8]{ctex}
\begin{document}

\section*{最终的确定性CD-SNDP MILP}
\begin{equation}
\begin{aligned}
(\mathbf{MILP})\quad \max_{v,f,p,x,z}\quad
&\sum_{(i,j)\in\mathcal A}-\frac{1}{\beta_c}\Biggl(
D_{ij}\lambda_{ij}+\sum_{h\in\mathcal H}U_{ij}^{h}z_{ij}^{h}
+\sum_{s\in\mathcal S}\sum_{k\in\mathcal K}\bar U_{ij}^{O}x_{ijsk}\\
&\hspace{42mm}+\beta_f\sum_{s\in\mathcal S}\sum_{k\in\mathcal K}\phi_{ijs}
\sum_{b=0}^{B_{sk}-1}2^b a_{ijskb}\Biggr)\\
&-\sum_{s\in\mathcal S}\sum_{k\in\mathcal K}c_{sk}^{\mathrm{FIX}}f_{sk}
-\sum_{(i,j)\in\mathcal A}\sum_{s\in\mathcal S}\sum_{k\in\mathcal K}c_{ijsk}^{\mathrm{VAR}}x_{ijsk}.
\end{aligned}\tag{19}
\end{equation}
在继承的运算符和需求约束的前提下：
\begin{align}
\sum_{s\in\mathcal S}v_{sk}&\leq V_k
&&\forall k\in\mathcal K,\tag{2}\\
f_{sk}&\leq W_{sk}v_{sk}
&&\forall s\in\mathcal S,\ \forall k\in\mathcal K,\tag{3}\\
\sum_{(i,j)\in\mathcal A}\delta_{ijl_s}x_{ijsk}&\leq Q_k f_{sk}
&&\forall l_s\in\mathcal L_s,\ \forall s\in\mathcal S,\ \forall k\in\mathcal K,\tag{4}\\
x_{ijsk}&\leq\sum_{l_s\in\mathcal L_s}\delta_{ijl_s}D_{ij}
&&\forall(i,j)\in\mathcal A,\ \forall s\in\mathcal S,\ \forall k\in\mathcal K,\tag{5}\\
p_{ij}&\geq0
&&\forall(i,j)\in\mathcal A,\tag{6}\\
v_{sk},f_{sk}&\in\mathbb N
&&\forall s\in\mathcal S,\ \forall k\in\mathcal K,\tag{7}\\
\sum_{s\in\mathcal S}\sum_{k\in\mathcal K}x_{ijsk}
+\sum_{h\in\mathcal H}z_{ij}^{h}&=D_{ij}
&&\forall(i,j)\in\mathcal A,\tag{9}\\
x_{ijsk}&\geq0
&&\forall(i,j)\in\mathcal A,\ \forall s\in\mathcal S,\ \forall k\in\mathcal K,\tag{10}\\
z_{ij}^{h}&\geq0
&&\forall(i,j)\in\mathcal A,\ \forall h\in\mathcal H.\tag{11}
\end{align}
最终的表述中增加了以下内容：频率扩展、乘积线性化，以及 KKT 析取行。
\begin{align}
f_{sk}&=\sum_{b=0}^{B_{sk}-1}2^b f_{skb}
&&\forall s\in\mathcal S,\ \forall k\in\mathcal K,\tag{20}\\
a_{ijskb}&\leq D_{ij}f_{skb}
&&\forall(i,j)\in\mathcal A,\ \forall s\in\mathcal S,\ \forall k\in\mathcal K,\ \forall b\in\mathcal B,\tag{21}\\
a_{ijskb}&\leq x_{ijsk}
&&\forall(i,j)\in\mathcal A,\ \forall s\in\mathcal S,\ \forall k\in\mathcal K,\ \forall b\in\mathcal B,\tag{22}\\
a_{ijskb}&\geq x_{ijsk}-D_{ij}(1-f_{skb})
&&\forall(i,j)\in\mathcal A,\ \forall s\in\mathcal S,\ \forall k\in\mathcal K,\ \forall b\in\mathcal B.\tag{23}
\end{align}
\begin{align}
\lambda_{ij}&\leq-\left(\bar U_{ij}^{O}+\beta_c p_{ij}
+\beta_f\sum_{s\in\mathcal S}\sum_{k\in\mathcal K}\phi_{ijs}f_{sk}\right)
&&\forall(i,j)\in\mathcal A,\tag{24}\\
-\left(\bar U_{ij}^{O}+\beta_c p_{ij}
+\beta_f\sum_{s\in\mathcal S}\sum_{k\in\mathcal K}\phi_{ijs}f_{sk}\right)
-\lambda_{ij}&\leq M_{ij}^{\mathrm I}y_{ij}^{\mathrm I}
&&\forall(i,j)\in\mathcal A,\tag{25}\\
\sum_{s\in\mathcal S}\sum_{k\in\mathcal K}x_{ijsk}&\leq M_{ij}^{\mathrm{II}}y_{ij}^{\mathrm{II}}
&&\forall(i,j)\in\mathcal A,\tag{26}\\
y_{ij}^{\mathrm I}+y_{ij}^{\mathrm{II}}&\leq1
&&\forall(i,j)\in\mathcal A,\tag{27}\\
\lambda_{ij}&\leq-U_{ij}^{h}
&&\forall(i,j)\in\mathcal A,\ \forall h\in\mathcal H,\tag{28}\\
-U_{ij}^{h}-\lambda_{ij}&\leq M_{ij}^{\mathrm I h}y_{ij}^{\mathrm I h}
&&\forall(i,j)\in\mathcal A,\ \forall h\in\mathcal H,\tag{29}\\
z_{ij}^{h}&\leq M_{ij}^{\mathrm{II}h}y_{ij}^{\mathrm{II}h}
&&\forall(i,j)\in\mathcal A,\ \forall h\in\mathcal H,\tag{30}\\
y_{ij}^{\mathrm I h}+y_{ij}^{\mathrm{II}h}&\leq1
&&\forall(i,j)\in\mathcal A,\ \forall h\in\mathcal H.\tag{31}
\end{align}
另外打印出来的变量域包括：
\begin{align}
f_{skb}&\in\{0,1\}
&&\forall s\in\mathcal S,\ \forall k\in\mathcal K,\ \forall b\in\mathcal B,\tag{32}\\
a_{ijskb}&\in\mathbb N
&&\forall(i,j)\in\mathcal A,\ \forall s\in\mathcal S,\ \forall k\in\mathcal K,\ \forall b\in\mathcal B,\tag{33}\\
y_{ij}^{\mathrm I},y_{ij}^{\mathrm{II}},y_{ij}^{\mathrm I h},y_{ij}^{\mathrm{II}h}&\in\{0,1\}
&&\forall(i,j)\in\mathcal A,\tag{34}\\
\lambda_{ij}&\geq0
&&\forall(i,j)\in\mathcal A.\tag{35}
\end{align}
所选的 CD-SNDP 还包含了后续的频率上限：
\begin{equation}
f_{sk}\leq35
\qquad\forall s\in\mathcal S,\ \forall k\in\mathcal K.
\tag{56}
\end{equation}

\section*{操作定义与变量角色}
\paragraph{网络与索引。}
有向运输网络是
\[
G=(\mathcal N,\mathcal A),\qquad
\mathcal A=\{(i,j):i,j\in\mathcal N,\ i\ne j\}.
\]
这里的\(\mathcal N\)表示终端集合，\(\mathcal K\)表示车辆类型集合，\(\mathcal S\)表示潜在的运输服务集合，而\(\mathcal H\)则代表竞争性的替代方案。对于每一项服务，\(s\)表示其路径序列，而\(\mathcal L_s\)则表示这些路径的索引，索引由\(l_s\)表示。每条路径都是一个弧线，且服务序列的起始点和终点都位于同一个节点上。服务过程中可能会有中途的停靠点。

\paragraph{固定操作数据。}
对于车辆类型\(k\)，\(V_k\)表示可用车辆的数量，而\(Q_k\)则表示这些车辆的集装箱处理能力，以TEU为单位。参数\(W_{sk}\)表示在规划时间范围内，一辆车辆类型\(k\)能够完成的最大服务周期数\(s\)。该论文将这一数值描述为最大运行时间除以周期时间，包括运输时间和在码头停留的时间；其具体计算过程见PDF第20页。
\[
W_{sk}=\frac{T^{\max}}{t_{sk}^{\mathrm{sail}}+t_{sk}^{\mathrm{port}}}.
\]
此处\(T^{\max}\)表示最大运营时间，\(t_{sk}^{\mathrm{sail}}\)表示航行时间，\(t_{sk}^{\mathrm{port}}\)表示在港口停留的时间。该表达式不进行四舍五入处理。

服务\(s\)的固定运营成本，其类型为\(k\)，为\(c_{sk}^{\mathrm{FIX}}\)欧元；使用该服务将货物从\(i\)运输到\(j\)时，每TEU的变动成本为\(c_{ijsk}^{\mathrm{VAR}}\)欧元。总需求为\(D_{ij}\)TEU。相关参数已定义。
\[
\delta_{ijl_s}=\begin{cases}
1,&\text{if a container traveling from }i\text{ to }j\text{ uses leg }l_s,\\
0,&\text{otherwise},
\end{cases}
\qquad (i,j)\in\mathcal A,\ s\in\mathcal S,\ l_s\in\mathcal L_s,
\]
\[
\phi_{ijs}=\begin{cases}
1,&\text{if both terminals }i\text{ and }j\text{ are contained in service }s,\\
0,&\text{otherwise},
\end{cases}
\qquad (i,j)\in\mathcal A,\ s\in\mathcal S.
\]

\paragraph{确定性效用定义。}
所选的通用实用工具规范，就是用于推导（19）的那个规范。
\begin{equation}
U_{ij}^{O}=\bar U_{ij}^{O}+\beta_c p_{ij}+\beta_f f_{ij}
=\bar U_{ij}^{O}+\beta_c p_{ij}
+\beta_f\sum_{s\in\mathcal S}\sum_{k\in\mathcal K}\phi_{ijs}f_{sk}.
\tag{17}
\end{equation}
特别是，
\[
f_{ij}=\sum_{s\in\mathcal S}\sum_{k\in\mathcal K}\phi_{ijs}f_{sk}.
\]
这些定义适用于OD对\((i,j)\in\mathcal A\)。量\(\bar U_{ij}^{O}\)是操作符效用的外生部分，而非全部效用\(U_{ij}^{O}\)。系数\(\beta_c\)和\(\beta_f\)分别用于权衡价格和频率。竞争性效用\(U_{ij}^{h}\)对于\((i,j)\in\mathcal A\)和\(h\in\mathcal H\)是固定数据。在这种通用的确定性表述中，它们仍然保持抽象状态，\(\bar U_{ij}^{O}\)也是如此；没有引入随机误差、抽样系数或实现指标。竞争是固定的且外生的，不会因操作符的决策而发生变化。

\paragraph{频率表示方式。}
二进制位数的数量在紧接着（19）之后的文字中有所说明：
\[
B_{sk}=\left\lceil\log_2(W_{sk}V_k+1)\right\rceil.
\]
在(19)和(20)中，求和的结果是从\(b=0\)到\(B_{sk}-1\)。变量\(f_{skb}\)是\(f_{sk}\)的二进制数字。辅助变量\(a_{ijskb}\)表示二进制数字\(f_{skb}\)与货物变量\(x_{ijsk}\)的乘积，这是通过打印的线性化公式(21)--(23)和域(33)来确定的，而不是通过额外的非线性等式来表示的。

\paragraph{决策变量。}
操作员选择了\(v_{sk}\)，即分配给服务\(s\)的、类型为\(k\)的车辆数量；\(f_{sk}\)表示使用这种类型车辆的频率；而\(p_{ij}\)则表示从\(i\)到\(j\)的每集装箱的价格。变量\(x_{ijsk}\)表示分配给该操作员服务的货物量，而\(z_{ij}^{h}\)则表示分配给竞争对手的替代方案中的货物量。原始表格中，货物量以TEU为单位表示，价格则以欧元/TEU为单位表示。

变量\(\lambda_{ij}\)被引入为与需求等式（9）相关的对偶变量。二元对\((y_{ij}^{\mathrm I},y_{ij}^{\mathrm{II}})\)和\((y_{ij}^{\mathrm I h},y_{ij}^{\mathrm{II}h})\)分别实现了操作符和竞争者互补性之间的分离。符号\(M_{ij}^{\mathrm I}\)、\(M_{ij}^{\mathrm{II}}\)、\(M_{ij}^{\mathrm I h}\)和\(M_{ij}^{\mathrm{II}h}\)是固定的大\(M\)系数；本文在此公式中未提供它们的公式或数值界限，仅以抽象形式保留它们。尽管在（19）中的打印最大化下标中仅列出了\(v,f,p,x,z\)，但频率数字、乘积变量、对偶变量和分离变量也是最终公式中的变量。

\paragraph{新增限制的范围。}
方程（56）被引入是因为，用于估算模式选择系数的数据中，每周不超过35项服务的IWT频率被包含在内。其打印范围是指每项服务的频率\(f_{sk}\)，而不是总出行频率\(f_{ij}\)。除了这个明确设定的上限外，案例研究中提到的车队数量、网络状况、效用系数以及运营时间等数值并未被强制执行。

\section*{来源资格保持不变，不进行修改。}
\begin{itemize}
\item PDF页面16上的最后一个块，明确继承了(2)--(7)和(9)--(11)。其中的“或”关系取代了之前的非线性互补行；这里既没有替换掉的那些行，也没有单独的低级优化问题。
\item 源代码中定义了标量数字的数量\(B_{sk}\)，但在(21)--(23)、(32)和(33)中使用了不同的集合符号\(\mathcal B\)，但没有明确定义该集合，也没有将其与服务相关的数字范围联系起来。打印出的总和以及相关的量词都得到了保留；而\(\mathcal B\)的定义则没有给出。
\item 方程（34）中包含了对竞标者进行索引的变量，但只打印了量词\(\forall(i,j)\in\mathcal A\)，而没有打印\(h\)量词。这种省略现象被保留了下来。在(7)和(33)中出现的符号\(\mathbb N\)也被保留下来，同时也没有添加关于零的任何约定。
\item 方程（33）使得乘积变量为整数，尽管前面的描述描述的是一种二进制连续乘积；而方程（10）仅对货物的非负性提出了要求。方程（35）则使得需求等式中的对偶变量为非负数，尽管等式乘数通常不受符号限制。这些打印出来的限制条件被保留下来，不会被放宽。
\item 在(19)中，频率项按照PDF页面15上印刷的相同服务产品顺序进行排列。相比之下，(17)则使用了每个运营商运输中的总OD频率。将这一总频率效用扩展到所有运营商的货物运输中，也会产生跨服务产品，而这些产品在(19)中并未出现。印刷版中的目标和效用定义保持不变，没有添加任何此类术语。
\item 通用公式将\(\beta_c\)作为除数进行除法运算，但并未规定其允许的范围或零系数的例外情况。也没有插入任何额外的符号或非零约束条件。在(17)之后的文字中，外部因素被标记为\(\bar U_{ij}\)，而公式中的上标\(O\)则被省略了；操作符号则遵循公式的排列顺序。
\end{itemize}

\paragraph{来源。}
Adrien Nicolet和Bilge Atasoy的论文《\emph{基于选择驱动的服务网络设计与定价，包括异质行为}》，arXiv:2311.15907v1，物理PDF页码为12–16和20。提供的页面图片展示了所选行在打印时的形式。该论文引用了Tawfik和Limbourg（2019）的论文作为参考资料，但其中关于转换细节的内容并未得到验证；也没有从该参考文献中引入任何约束或约定。

\end{document}
